\subsection{Right band acts and the seven-state double-band gate}
\begin{lemma}[Proved: complete right-$P_5$ act normal forms]
\label{CP2:A-P5-normal-forms}
Let $P_5=\{1,e_{00},e_{01},e_{10},e_{11}\}$, with
$e_{ij}e_{kl}=e_{il}$. Every finite right $P_5$-set is determined,
up to an equivariant bijection, by $a$ singleton core orbits,
$b$ two-element core orbits with labelled columns, and a multiset
of $t$ ordered pairs of these $a+b$ orbits. Its cardinality is
$a+2b+t$. The only ambiguity is permutation of the singleton
orbits, permutation of the double orbits, and permutation of
equal-type transient points. Consequently the numbers on
$1,\ldots,7$ points are $1,3,6,17,44,130,383$.
\end{lemma}
\begin{proof}
For a right act $X$, put $Y=X\{e_{ij}\}$. If $y=xe_{ij}$,
then $ye_{kl}=xe_{il}$ is independent of $k$. Thus its orbit
consists of one point or of two points canonically labelled
by the column $l$. Two such orbits meeting at a point coincide,
since applying either column operation to that point recovers
both labelled outputs. These are precisely the core orbits.
For $x\in X\setminus Y$, each row $i$ selects the orbit
$\{xe_{i0},xe_{i1}\}$, and the column selects its labelled
point. These choices are unrestricted and satisfy every act
equation. Transient points with the same ordered orbit pair
are interchangeable; core orbit sizes are preserved. This
proves both existence and completeness of the normal form.

The producer \path{code/p5_right_acts.py} enumerates every
$a,b,t$ with $a+2b+t=s$, all multisets of ordered orbit pairs,
and their $S_a\times S_b$ orbits. An independent count uses
Burnside's lemma. For a permutation of the orbit set, let its
action on ordered pairs have cycle lengths $d_1,\ldots,d_r$.
The fixed multisets of size $t$ are counted by the coefficient
of $x^t$ in $\prod_i(1-x^{d_i})^{-1}$. The separate program
\path{code/check_p5_normal_forms_burnside.py} averages these
integer coefficients and checks every parameter case against
the literal normal forms, as well as every action equation.
It obtains the stated counts. The input tables, all parameter
counts and digests are under \path{outputs/p5_bisets/}.
\end{proof}

\begin{proposition}[Proved: the $P_5$--$P_5$ seven-state case]
\label{CP2:A-P5-P5-seven}
No category with two objects, both corner monoids isomorphic
to $P_5$, seven arrows in one direction between the objects,
and no arrows in the reverse direction has type $\CP^2$.
The total arrow count, including identities, is seventeen.
\end{proposition}
\begin{proof}
Such a category is exactly a commuting left and right $P_5$
action on seven states. Fix a representative of the right act
from Lemma~\ref{CP2:A-P5-normal-forms}. The transpose of the
band indices identifies $P_5^{\mathrm{op}}$ with $P_5$.
Taking all state relabellings of all $383$ normal forms therefore
gives every labelled left action. The new program
\path{code/p5_bisets.cpp} constructs exactly $1038598$ such
left actions. It verifies that the normal-form state orbits are
disjoint and that their cardinality is
$\sum 7!/|\operatorname{Aut}(X)|$.

For each fixed right act it computes every transformation of
the seven states commuting with its four band operations.
It tries every labelled left action whose four maps are in
that centralizer. The automorphism group of the right act
acts on these choices; taking the least member of each orbit
loses only an equivariant state relabelling. Its full group
is obtained by testing every one of the $7!$ state permutations.
The program separately verifies that the number of compatible
labelled actions equals the sum of these orbit cardinalities.
In the trivial-right-action case the already complete $383$
left normal forms are used, with the same orbit-weight check.
No monoid automorphism quotient, opposite-category quotient,
retraction cut or conditional corner reduction is used.

There are $1904633$ compatible labelled left actions over the
right-act representatives and $56597$ resulting state orbits
of bisets. Exactly $15606$ fail $H_1(-;\mathbb F_2)=0$;
another $24704$ fail $\dim H_2(-;\mathbb F_2)=1$.
Every remaining one of the $16287$ literal categories has a
nonzero two-cocycle whose square is the coboundary of a supplied
three-cochain. The independent program
\path{code/check_p5_bisets.cpp}, which imports no producer
code, reconstructs all typed tables and normalized chains.
It checks $12622425$ associativity equations, $7557168$
cocycle equations and $37525248$ square-primitive equations.
All records pass. Since the nonzero class of $H^2(\CP^2;
\mathbb F_2)$ has nonzero square, every case is excluded.
The exact scope, complete literal stream, digests and checker
receipt are in \path{outputs/p5_bisets_live/} in the checkpoint.
This proposition asserts only the double-$P_5$ case; mixed
order-five corner monoids remain a separate gate.
\end{proof}

\begin{corollary}[Proved: the double-$P_5$ cases through seventeen arrows]
\label{CP2:A-P5-P5-through17}
A two-object category with both corners $P_5$, arrows only in one
cross direction, and at most seventeen total arrows cannot have
classifying space $\mathbb{CP}^2$.
\end{corollary}
\begin{proof}
The preceding construction and independent literal check apply unchanged
with $s=1,\ldots,6$ states. The numbers of state-isomorphism orbits are,
respectively, $1,6,32,181,1124,7667$. After the two mod-two homology
filters, the retained counts are $0,4,10,60,360,2320$. Every retained
category has a supplied and independently checked square primitive.
The zero-state case is disconnected. Exact streams and independent
receipts are in \path{outputs/p5_bisets_small/}; the separate checker
is \path{code/check_p5_bisets_all_sizes.cpp}. Together with the
seven-state proposition this covers every possible cross size.
\end{proof}
